\documentclass[10pt,a3paper,landscape]{article}
\usepackage[margin=6mm]{geometry}
\usepackage[T1]{fontenc}
\usepackage[utf8]{inputenc}
\usepackage[french]{babel}
\usepackage{lmodern}
\usepackage{microtype}
\makeatletter\def\input@path{{../../../Style/}}\makeatother
\NeedsTeXFormat{LaTeX2e}
\ProvidesPackage{TLPosterCarteMentale}[2026/08/03 Posters et cartes mentales SI]

\RequirePackage{tikz}
\RequirePackage[most]{tcolorbox}
\RequirePackage{xcolor}
\RequirePackage{amsmath,amssymb}
\RequirePackage{graphicx}
\usetikzlibrary{positioning,arrows.meta,calc,shapes.geometric,mindmap,backgrounds,fit}

\definecolor{TLPosterBlue}{RGB}{24,86,141}
\definecolor{TLPosterLight}{RGB}{234,243,250}
\definecolor{TLPosterAccent}{RGB}{232,126,34}
\definecolor{TLPosterGreen}{RGB}{52,152,102}
\definecolor{TLPosterRed}{RGB}{192,57,43}
\definecolor{TLPosterGray}{RGB}{80,90,100}

\newcommand{\TLPosterTitre}[2]{%
  \begin{tcolorbox}[enhanced,boxrule=0pt,arc=0pt,colback=TLPosterBlue,left=6mm,right=6mm,top=4mm,bottom=4mm]
    {\color{white}\bfseries\LARGE #1}\hfill{\color{white}\large #2}
  \end{tcolorbox}%
}

\newtcolorbox{TLPosterBloc}[2][]{
  enhanced,breakable,
  colback=TLPosterLight,
  colframe=TLPosterBlue,
  colbacktitle=TLPosterBlue,
  coltitle=white,
  fonttitle=\bfseries,
  title={#2},
  boxrule=0.8pt,
  arc=2mm,
  left=3mm,right=3mm,top=2mm,bottom=2mm,
  #1
}

\newtcolorbox{TLPosterFormule}[1][]{
  enhanced,
  colback=white,
  colframe=TLPosterAccent,
  boxrule=1pt,
  arc=2mm,
  fontupper=\bfseries,
  halign=center,
  #1
}

\newcommand{\TLPosterMethode}[1]{%
  \begin{tcolorbox}[enhanced,colback=TLPosterGreen!8,colframe=TLPosterGreen,title=Méthode,fonttitle=\bfseries]
  #1
  \end{tcolorbox}%
}

\newcommand{\TLPosterAttention}[1]{%
  \begin{tcolorbox}[enhanced,colback=TLPosterRed!6,colframe=TLPosterRed,title=Point de vigilance,fonttitle=\bfseries]
  #1
  \end{tcolorbox}%
}

\tikzset{
  TLmindroot/.style={concept,concept color=TLPosterBlue,text=white,font=\bfseries\large,minimum size=34mm},
  TLmindlevel1/.style={level 1 concept/.append style={font=\bfseries,minimum size=23mm,sibling angle=45}},
  TLmindlevel2/.style={level 2 concept/.append style={font=\small,minimum size=17mm}},
  TLarrow/.style={-{Stealth[length=3mm]},thick,draw=TLPosterBlue},
  TLbox/.style={rounded corners=2mm,draw=TLPosterBlue,fill=TLPosterLight,align=center,inner sep=3mm}
}

\newenvironment{TLCarteMentale}[1]{%
  \begin{tikzpicture}[mindmap,grow cyclic,concept color=TLPosterBlue,every node/.style={concept},TLmindlevel1,TLmindlevel2]
  \node[TLmindroot]{#1}
}{;
  \end{tikzpicture}
}

\newcommand{\TLBranche}[3][]{child[concept color=#2]{node[#1]{#3}}}

\endinput

\pagestyle{empty}
\renewcommand{\familydefault}{\sfdefault}
\begin{document}
\begin{tikzpicture}
\TLMapBase{Électronique de puissance}{Choisir le cas $\rightarrow$ séquencer $\rightarrow$ tracer $\rightarrow$ calculer $\rightarrow$ vérifier}
\TLMapBranchTL{À connaître}{\textbullet\ source tension : ne jamais court-circuiter ; source courant : ne jamais ouvrir\\[1mm]\textbullet\ $L$ impose continuité du courant ; $C$ impose continuité de la tension\\[1mm]\textbullet\ passant : $u_K=0$ ; bloqué : $i_K=0$\\[1mm]\textbullet\ $\alpha=t_{on}/T$, $\langle u_L\rangle=0$ en régime périodique\\[1mm]\textbullet\ $P=\langle ui\rangle$, $S=UI$, $\lambda=P/S$}
\TLMapBranchTR{Choisir le convertisseur}{\textbf{Buck :} $U_s<U_e$\\[1mm]\textbf{Boost :} $U_s>U_e$\\[1mm]\textbf{2 quadrants :} une réversibilité\\[1mm]\textbf{4 quadrants :} tension et courant réversibles\\[1mm]\textbf{PD2 :} AC $\rightarrow$ DC non commandé\\[1mm]\textbf{PFC :} courant réseau sinusoïdal et en phase}
\TLMapBranchML{Méthode hacheur}{1. Repérer $T$ et les intervalles.\\[1mm]2. Identifier passants/bloqués.\\[1mm]3. Remplacer par CC/CO et redessiner.\\[1mm]4. Écrire les lois, surtout $u_L=L\,di/dt$.\\[1mm]5. Tracer les chronogrammes.\\[1mm]6. Utiliser $\langle u_L\rangle=0$.\\[1mm]7. Calculer ondulation et contraintes.}
\TLMapBranchMR{Méthode AC / puissance}{\textbf{PD2 :} choisir la diagonale conductrice sur chaque alternance ; $u_s=|u_e|$ ; charge inductive $\Rightarrow$ courant continu.\\[1mm]\textbf{Puissance :} partir de $p(t)=ui$ ; $P=UI\cos\varphi$ seulement en sinusoïdal.\\[1mm]\textbf{PFC :} $i_{ref}\propto|u_e|$ ; boucle courant rapide, boucle tension lente.}
\TLMapBranchBL{Composants / pertes}{\textbullet\ diode : commutation spontanée\\[1mm]\textbullet\ MOSFET / IGBT : amorçage et blocage commandés\\[1mm]\textbullet\ transistor + diode antiparallèle : réversible en courant\\[1mm]\textbullet\ $P_{com}=f_s(E_{on}+E_{off})$\\[1mm]\textbullet\ vérifier $U_{max}$, $I_{max}$, $I_{eff}$, pertes, $T_j$}
\TLMapBranchBR{Vérifications}{\textbullet\ associations de sources autorisées ?\\[1mm]\textbullet\ jamais deux interrupteurs d'un bras passants simultanément\\[1mm]\textbullet\ continuité de $i_L$ et de $u_C$ respectée ?\\[1mm]\textbullet\ signes et réversibilités cohérents ?\\[1mm]\textbullet\ moyenne, efficace, ondulation, rendement plausibles ?\\[1mm]\textbullet\ régime déformé : $S^2=P^2+Q^2+D^2$}
\TLMapRibbon{Identifier le cas avant de calculer : la méthode n'est pas la même pour un hacheur, un redresseur ou une étude de qualité d'énergie.}
\end{tikzpicture}
\end{document}
